\text{$i = j$ and $(i,i+1)$, $(i+1,i+1) \in S(\mu) \setminus D$.}
\end{gather*}
\item\label{exam3.19_b} %[(b)]
If $u = (i,j) \in D$ is $D$-active, then we define an ordinary and $K$-theoretical 
elementary excitation by putting
\begin{align*}
\alpha_u(D) &= D \setminus \{ (i,j) \} \cup \{ (i+1,j+1) \},
\\
\alpha^*_u(D) &= D \cup \{ (i+1,j+1) \},
\end{align*}
respectively.
\item\label{exam3.19_c} %[(c)]
If $u = (i,j) \in D$ is $D$-active, then we define
$$
B(\alpha_u(D))
 =
\begin{cases}
 B(D) \setminus \{ (i,j+1), (i+1,j) \} \cup \{ (i,j) \} 
 &\text{if $i<j$,}
\\
 B(D) \setminus \{ (i,j+1) \} \cup \{ (i,j) \} 
 &\text{if $i=j$.}
\end{cases}
$$
\end{enumerate}
This notion of excited diagrams is the same as Ikeda--Naruse's excited diagrams of type $I$ 
introduced in~\cite{IN1}.

For example, if $P = S(5,4,2,1)$ and $F = S(3,1)$, 
then there are $10$ excited diagrams of $F$ in $P$ as a heap for the type $B$ Weyl group.
See Figure~\ref{fig:shifted-B_excited}.
\begin{figure}[!htb]
$$
\setlength{\unitlength}{1.1pt}
\begin{CD}
\raisebox{-20pt}{
\begin{picture}(50,40)
\fboxsep=0mm
\put(0,30){\colorbox[gray]{0.7}{\makebox(10,10){}}}
\put(10,30){\colorbox[gray]{0.7}{\makebox(10,10){}}}
\put(20,30){\colorbox[gray]{0.7}{\makebox(10,10){}}}
\put(10,20){\colorbox[gray]{0.7}{\makebox(10,10){}}}
\put(0,40){\line(1,0){50}}
\put(0,30){\line(1,0){50}}
\put(10,20){\line(1,0){40}}
\put(20,10){\line(1,0){20}}
\put(30,0){\line(1,0){10}}
\put(0,30){\line(0,1){10}}
\put(10,20){\line(0,1){20}}
\put(20,10){\line(0,1){30}}
\put(30,0){\line(0,1){40}}
\put(40,0){\line(0,1){40}}
\put(50,20){\line(0,1){20}}
\end{picture}
}
@>>>
\raisebox{-20pt}{
\begin{picture}(50,40)
\fboxsep=0mm
\put(0,30){\colorbox[gray]{0.7}{\makebox(10,10){}}}
\put(10,30){\colorbox[gray]{0.7}{\makebox(10,10){}}}
\put(30,20){\colorbox[gray]{0.7}{\makebox(10,10){}}}
\put(10,20){\colorbox[gray]{0.7}{\makebox(10,10){}}}
\put(0,40){\line(1,0){50}}
\put(0,30){\line(1,0){50}}
\put(10,20){\line(1,0){40}}
\put(20,10){\line(1,0){20}}
\put(30,0){\line(1,0){10}}
\put(0,30){\line(0,1){10}}
\put(10,20){\line(0,1){20}}
\put(20,10){\line(0,1){30}}
\put(30,0){\line(0,1){40}}
\put(40,0){\line(0,1){40}}
\put(50,20){\line(0,1){20}}
\put(20,30){\makebox(10,10){$\times$}}
\end{picture}
}
\\[5pt]
@VVV
@VVV
\\[5pt]
\raisebox{-20pt}{
\begin{picture}(50,40)
\fboxsep=0mm
\put(0,30){\colorbox[gray]{0.7}{\makebox(10,10){}}}
\put(10,30){\colorbox[gray]{0.7}{\makebox(10,10){}}}
\put(20,30){\colorbox[gray]{0.7}{\makebox(10,10){}}}
\put(20,10){\colorbox[gray]{0.7}{\makebox(10,10){}}}
\put(0,40){\line(1,0){50}}
\put(0,30){\line(1,0){50}}
\put(10,20){\line(1,0){40}}
\put(20,10){\line(1,0){20}}
\put(30,0){\line(1,0){10}}
\put(0,30){\line(0,1){10}}
\put(10,20){\line(0,1){20}}
\put(20,10){\line(0,1){30}}
\put(30,0){\line(0,1){40}}
\put(40,0){\line(0,1){40}}
\put(50,20){\line(0,1){20}}
\put(10,20){\makebox(10,10){$\times$}}
\end{picture}
}
@>>>
\raisebox{-20pt}{
\begin{picture}(50,40)
\fboxsep=0mm
\put(0,30){\colorbox[gray]{0.7}{\makebox(10,10){}}}
\put(10,30){\colorbox[gray]{0.7}{\makebox(10,10){}}}
\put(30,20){\colorbox[gray]{0.7}{\makebox(10,10){}}}
\put(20,10){\colorbox[gray]{0.7}{\makebox(10,10){}}}
\put(0,40){\line(1,0){50}}
\put(0,30){\line(1,0){50}}
\put(10,20){\line(1,0){40}}
\put(20,10){\line(1,0){20}}
\put(30,0){\line(1,0){10}}
\put(0,30){\line(0,1){10}}
\put(10,20){\line(0,1){20}}
\put(20,10){\line(0,1){30}}
\put(30,0){\line(0,1){40}}
\put(40,0){\line(0,1){40}}
\put(50,20){\line(0,1){20}}
\put(20,30){\makebox(10,10){$\times$}}
\put(10,20){\makebox(10,10){$\times$}}
\end{picture}
}
@>>>
\raisebox{-20pt}{
\begin{picture}(50,40)
\fboxsep=0mm
\put(0,30){\colorbox[gray]{0.7}{\makebox(10,10){}}}
\put(20,20){\colorbox[gray]{0.7}{\makebox(10,10){}}}
\put(20,10){\colorbox[gray]{0.7}{\makebox(10,10){}}}
\put(30,20){\colorbox[gray]{0.7}{\makebox(10,10){}}}
\put(0,40){\line(1,0){50}}
\put(0,30){\line(1,0){50}}
\put(10,20){\line(1,0){40}}
\put(20,10){\line(1,0){20}}
\put(30,0){\line(1,0){10}}
\put(0,30){\line(0,1){10}}
\put(10,20){\line(0,1){20}}
\put(20,10){\line(0,1){30}}
\put(30,0){\line(0,1){40}}
\put(40,0){\line(0,1){40}}
\put(50,20){\line(0,1){20}}
\put(10,30){\makebox(10,10){$\times$}}
\end{picture}
}
@>>>
\raisebox{-20pt}{
\begin{picture}(50,40)
\fboxsep=0mm
\put(10,20){\colorbox[gray]{0.7}{\makebox(10,10){}}}
\put(20,20){\colorbox[gray]{0.7}{\makebox(10,10){}}}
\put(30,20){\colorbox[gray]{0.7}{\makebox(10,10){}}}
\put(20,10){\colorbox[gray]{0.7}{\makebox(10,10){}}}
\put(0,40){\line(1,0){50}}
\put(0,30){\line(1,0){50}}
\put(10,20){\line(1,0){40}}
\put(20,10){\line(1,0){20}}
\put(30,0){\line(1,0){10}}
\put(0,30){\line(0,1){10}}
\put(10,20){\line(0,1){20}}
\put(20,10){\line(0,1){30}}
\put(30,0){\line(0,1){40}}
\put(40,0){\line(0,1){40}}
\put(50,20){\line(0,1){20}}
\put(0,30){\makebox(10,10){$\times$}}
\end{picture}
}
\\[5pt]
@VVV
@VVV
@VVV
@VVV
\\[5pt]
\raisebox{-20pt}{
\begin{picture}(50,40)
\fboxsep=0mm
\put(0,30){\colorbox[gray]{0.7}{\makebox(10,10){}}}
\put(10,30){\colorbox[gray]{0.7}{\makebox(10,10){}}}
\put(20,30){\colorbox[gray]{0.7}{\makebox(10,10){}}}
\put(30,0){\colorbox[gray]{0.7}{\makebox(10,10){}}}
\put(0,40){\line(1,0){50}}
\put(0,30){\line(1,0){50}}
\put(10,20){\line(1,0){40}}
\put(20,10){\line(1,0){20}}
\put(30,0){\line(1,0){10}}
\put(0,30){\line(0,1){10}}
\put(10,20){\line(0,1){20}}
\put(20,10){\line(0,1){30}}
\put(30,0){\line(0,1){40}}
\put(40,0){\line(0,1){40}}
\put(50,20){\line(0,1){20}}
\put(10,20){\makebox(10,10){$\times$}}
\put(20,10){\makebox(10,10){$\times$}}
\end{picture}
}
@>>>
\raisebox{-20pt}{
\begin{picture}(50,40)
\fboxsep=0mm
\put(0,30){\colorbox[gray]{0.7}{\makebox(10,10){}}}
\put(10,30){\colorbox[gray]{0.7}{\makebox(10,10){}}}
\put(30,20){\colorbox[gray]{0.7}{\makebox(10,10){}}}
\put(30,0){\colorbox[gray]{0.7}{\makebox(10,10){}}}
\put(0,40){\line(1,0){50}}
\put(0,30){\line(1,0){50}}
\put(10,20){\line(1,0){40}}
\put(20,10){\line(1,0){20}}
\put(30,0){\line(1,0){10}}
\put(0,30){\line(0,1){10}}
\put(10,20){\line(0,1){20}}
\put(20,10){\line(0,1){30}}
\put(30,0){\line(0,1){40}}
\put(40,0){\line(0,1){40}}
\put(50,20){\line(0,1){20}}
\put(20,30){\makebox(10,10){$\times$}}
\put(10,20){\makebox(10,10){$\times$}}
\put(20,10){\makebox(10,10){$\times$}}
\end{picture}
}
@>>>
\raisebox{-20pt}{
\begin{picture}(50,40)
\fboxsep=0mm
\put(0,30){\colorbox[gray]{0.7}{\makebox(10,10){}}}
\put(20,20){\colorbox[gray]{0.7}{\makebox(10,10){}}}
\put(30,20){\colorbox[gray]{0.7}{\makebox(10,10){}}}
\put(30,0){\colorbox[gray]{0.7}{\makebox(10,10){}}}
\put(0,40){\line(1,0){50}}
\put(0,30){\line(1,0){50}}
\put(10,20){\line(1,0){40}}
\put(20,10){\line(1,0){20}}
\put(30,0){\line(1,0){10}}
\put(0,30){\line(0,1){10}}
\put(10,20){\line(0,1){20}}
\put(20,10){\line(0,1){30}}
\put(30,0){\line(0,1){40}}
\put(40,0){\line(0,1){40}}
\put(50,20){\line(0,1){20}}
\put(10,30){\makebox(10,10){$\times$}}
\put(20,10){\makebox(10,10){$\times$}}
\end{picture}
}
@>>>
\raisebox{-20pt}{
\begin{picture}(50,40)
\fboxsep=0mm
\put(10,20){\colorbox[gray]{0.7}{\makebox(10,10){}}}
\put(20,20){\colorbox[gray]{0.7}{\makebox(10,10){}}}
\put(30,20){\colorbox[gray]{0.7}{\makebox(10,10){}}}
\put(30,0){\colorbox[gray]{0.7}{\makebox(10,10){}}}
\put(0,40){\line(1,0){50}}
\put(0,30){\line(1,0){50}}
\put(10,20){\line(1,0){40}}
\put(20,10){\line(1,0){20}}
\put(30,0){\line(1,0){10}}
\put(0,30){\line(0,1){10}}
\put(10,20){\line(0,1){20}}
\put(20,10){\line(0,1){30}}
\put(30,0){\line(0,1){40}}
\put(40,0){\line(0,1){40}}
\put(50,20){\line(0,1){20}}
\put(0,30){\makebox(10,10){$\times$}}
\put(20,10){\makebox(10,10){$\times$}}
\end{picture}
}
\end{CD}
$$
\caption{Excited diagrams in $S(5,4,2,1)$ viewed as a type $B$ heap}
\label{fig:shifted-B_excited}
\end{figure}
\end{exam}

\section{%
Equivariant \texorpdfstring{$K$}{K}-theory of Kac--Moody partial flag varieties
}\label{sect4}

In this section we review the basic properties of the equivariant $K$-theory 
of thick flag varieties following~\cite[Section~3]{LSS}, 
and rephrase the Billey-type formula and the Chevalley-type formula 
in terms of combinatorics of $d$-complete posets.

\subsection{Equivariant \texorpdfstring{$K$}{K}-theory and localization}\label{subsec4.1}

Let $A = (a_{ij})_{i,j \in I}$ be a symmetrizable generalized Cartan matrix, 
and $\Gamma$ the corresponding Dynkin diagram with node set $I$.
Then the associated Kac--Moody group $\mathcal{G}$ over $\Comp$ is constructed from the following data:
the weight lattice ($\Int$-module) $\Lambda$, the simple roots $\Pi = \{ \alpha_i : i \in I \}$, 
the simple coroots $\Pi^\vee = \{ \alpha^\vee_i :i \in I \}$, 
and the fundamental weights $\{ \lambda_i : i \in I \}$ (see the beginning of Subsection~2.2).

In what follows, we fix a subset $J$ of $I$.
Let $\BB$ be a Borel subgroup corresponding to the positive system $\Phi_+$ 
and $\TT \subset \BB$ a maximal torus.
Let $\PP_-$ be the opposite parabolic subgroup corresponding to the subset $J$, 
which contains the opposite Borel subgroup $\BB_-$ such that $\BB \cap \BB_- = \TT$.
Then we can introduce the Kashiwara thick partial flag variety $\XX = \GG/\PP_-$.
(We refer the readers to~\cite{Ka} for a construction of $\XX$.)

Let $W_J$ be the parabolic subgroup of $W$ corresponding to $J$ 
and $W^J$ be the set of minimum length coset representatives of $W/W_J$.
For each element $v \in W^J$, we put $\XX^\circ_v = \BB v \PP_-/\PP_-$ 
and $\XX_v = \overline{\XX^\circ_v}$, the Zariski closure of $\XX^\circ_v$, 
which are called the Schubert cell and the Schubert variety respectively.
Then $\XX_v$ has codimension $l(v)$ in $\XX$ and
$$
\XX_v = \bigsqcup_{w \in W^J, \, w \ge v} \XX^\circ_w.
$$

Let $K_\TT(\XX)$ be the $\TT$-equivariant $K$-theory of $\XX$.
Then $K_\TT(\XX)$ has a commutative associative $K_\TT(\point)$-algebra structure.
Here the $\TT$-equivalent $K$-theory $K_\TT(\point)$ of a point is isomorphic 
to the group algebra $\Int[\Lambda]$ with basis $\{ e^\lambda : \lambda \in \Lambda \}$, 
and to the representation ring $R(\TT)$ of $\TT$.
For each $v \in W^J$, let $[\OO_v]$ be the class of the structure sheaf $\OO_v$ of $\XX_v$ 
in $K_\TT(\XX)$ and call it the equivariant Schubert class.
Then we have
$$
K_\TT(\XX) \cong \prod_{v \in W^J} K_\TT(\point) [\OO_v].
$$
Any elements of $K_\TT(\XX)$ is a (possibly infinite) $K_\TT(\point)$-linear combination 
of the equivariant Schubert classes.

Each $w \in W^J$ gives a $\TT$-fixed point $e_w = w \PP_-/\PP_- \in \XX$.
Then the inclusion map $\iota_w : \{ e_w \} \to \XX$ induces the pull-back ring homomorphism, 
called the localization map at $w$,
$$
\iota_w^* : K_{\TT}(\XX) \to K_{\TT}(e_w) \cong \Int[\Lambda].
$$
If $\LL^\lambda$ is the line bundle on $\XX$ corresponding to a weight $\lambda \in \Lambda$, 
then the image of the class $[\LL^\lambda]$ under the localization map is given by
$\iota_w^*( [\LL^\lambda] ) = e^{w \lambda}$.
For two elements $v$, $w \in W^J$, we denote by $\xi^v|_w$ 
the image of the $\TT$-equivariant Schubert class $\xi^v = [\OO_v] \in K_{\TT}(\XX)$ 
under the localization map $\iota_w^*$:
$$
\xi^v|_w = \iota_w^* ( [\OO_v] ).
$$
Then the Billey-type formula for the equivariant $K$-theory can be stated as follows:

\setbox\toto=\hbox {\cite[Proposition~2.10]{LSS}}
\begin{prop}[\box\toto]
\label{prop:Billey}
%\textup{(\cite[Proposition~2.10]{LSS})}
Let $v$, $w \in W^J$, and fix a reduced expression $w = s_{i_1} s_{i_2} \dots s_{i_N}$ of $w$.
Then we have
\begin{equation}
\label{eq:Billey}
\xi^v|_w 
 =
\sum_{(k_1, \dots, k_r)}
 (-1)^{r-l(v)} \prod_{a=1}^r \left( 1 - e^{\beta^{(k_a)}} \right),
\end{equation}
where the summation is taken over all sequences $(k_1, \dots, k_r)$ 
such that $1 \le k_1 < k_2 < \dots < k_r \le N$ and $s_{i_{k_1}} * \dots * s_{i_{k_r}} = v$ 
(with respect to the Demazure product), 
and $\beta^{(k)}$ is given by $\beta^{(k)} = s_{i_1} \dots s_{i_{k-1}}(\alpha_{i_k})$ for $1 \le k \le N$.
\end{prop}

By using Lemma~\ref{lem:red}~\ref{lemma3.12_a}, we can deduce the following corollary from Proposition~\ref{prop:Billey}.

\begin{coro} \label{cor:Billey}\ 

\begin{enumerate}[(a)]
\item\label{coro4.2_a} %[(a)]
For $w \in W^J$, we have
\begin{equation}
\label{eq:Billey2}
\xi^w|_w = \prod_{k=1}^N \left( 1 - e^{\beta^{(k)}} \right).
\end{equation}
In particular, $\xi^w|_w \neq 0$.
\item\label{coro4.2_b} %[(b)]
Let $v$, $w \in W^J$.
If $\xi^v|_w \neq 0$, then we have $v \le w$ in the Bruhat order.
\end{enumerate}
\end{coro}

\subsection{%
\texorpdfstring{E$\mkern -0.5mu$q$\mkern -0.5mu$u$\mkern -0.5mu$i$\mkern -0.5mu$v$\mkern -0.5mu$a$\mkern -0.5mu$r$\mkern -0.5mu$i$\mkern -0.5mu$a$\mkern -0.5mu$n$\mkern -0.5mu$t $\mkern -0.5muK\mkern -0.5mu$-$\mkern -0.5mu$t$\mkern -0.5mu$h$\mkern -0.5mu$e$\mkern -0.5mu$o$\mkern -0.5mu$r$\mkern -0.5mu$e$\mkern -0.5mu$t$\mkern -0.5mu$i$\mkern -0.5mu$c$\mkern -0.5mu$a$\mkern -0.5mu$l$\mkern -0.5mu$ L$\mkern -0.5mu$i$\mkern -0.5mu$t$\mkern -0.5mu$t$\mkern -0.5mu$l$\mkern -0.5mu$e$\mkern -0.5mu$w$\mkern -0.5mu$o$\mkern -0.5mu$o$\mkern -0.5mu$d$\mkern -0.5mu$--$\mkern -0.5mu$R$\mkern -0.5mu$i$\mkern -0.5mu$c$\mkern -0.5mu$h$\mkern -0.5mu$a$\mkern -0.5mu$r$\mkern -0.5mu$d$\mkern -0.5mu$s$\mkern -0.5mu$o$\mkern -0.5mu$n$\mkern -0.5mu$ c$\mkern -0.5mu$o$\mkern -0.5mu$e$\mkern -0.5mu$f$\mkern -0.5mu$f$\mkern -0.5mu$i$\mkern -0.5mu$c$\mkern -0.5mu$i$\mkern -0.5mu$e$\mkern -0.5mu$n$\mkern -0.5mu$t$\mkern -0.5mu$s}{Equivariant K-theoretical Littlewood--Richardson coefficients}
}\label{subsec4.2} 

We consider the structure constants for the multiplication in $K_{\TT}(\XX)$ 
with respect to the equivariant Schubert classes.
For $u$, $v$, $w \in W^J$, we denote by $c^w_{u,v} \in K_{\TT}(\point)$ 
the structure constant determined by
$$
[\OO_u] [\OO_v] = \sum_{w \in W^J} c^w_{u,v} [\OO_w].
$$

\begin{lemm}
\label{lem:LR}
If $c^w_{u,v} \neq 0$, then $u \le w$ and $v \le w$.
\end{lemm}

\begin{proof}
We use the induction on $l(w)$ to prove that, 
if $u \not\le w$ or $v \not\le w$, then $c^w_{u,v} = 0$.
Assume that $u \not\le w$ or $v \not\le w$.
By apply the localization map $\iota_w^*$ to 
$[\OO_u] [\OO_v] = \sum_{x \in W^J} c^x_{u,v} [\OO_x]$ 
and then by using Corollary~\ref{cor:Billey}~\ref{coro4.2_b}, we have
$$
\left( \xi^u|_w \right) \cdot \left( \xi^v|_w \right)
 =
\sum_{x \le w} c^x_{u,v} \xi^x|_w.
$$
If there exists an element $x \in W^J$ satisfying $x < w$ and $c^x_{u,v} \neq 0$, 
then we have $u \le x$ and $v \le x$ by the induction hypothesis, 
and hence $u \le w$ and $v \le w$, which contradicts to the assumption.
Hence, by using Corollary~\ref{cor:Billey}~\ref{coro4.2_b} and the assumption, we have
$$
0 = c^w_{u,v} \xi^w|_w.
$$
Since $\xi^w_w \neq 0$ (Corollary~\ref{cor:Billey}~\ref{coro4.2_a}), we obtain $c^w_{u,v} = 0$.
\end{proof}

\begin{prop}
\label{prop:LR=xi}
For $v$, $w \in W^J$, we have
\begin{equation}
\label{eq:LR=xi}
c^w_{v,w} = \xi^v|_w.
\end{equation}
\end{prop}

\begin{proof}
By applying the localization map $\iota_w^*$ to 
$[\OO_v] [\OO_w] = \sum_{x \in W^J} c^x_{v,w} [\OO_x]$ 
and then by using Corollary~\ref{cor:Billey}~\ref{coro4.2_b}, we have
$$
\left( \xi^v|_w \right) \cdot \left( \xi^w|_w \right)
 =
\sum_{x \le w} c^x_{v,w} \xi^x|_w.
$$
By Lemma~\ref{lem:LR}, we see that $c^x_{v,w} = 0$ unless $w \le x$.
By Corollary~\ref{cor:Billey}~\ref{coro4.2_b}, we see that $\xi^x|_w = 0$ unless $x \le w$.
Hence we have
$$
\left( \xi^v|_w \right) \cdot \left( \xi^w|_w \right)
 =
c^w_{v,w} \xi^w|_w.
$$
Since $\xi^w|_w \neq 0$ (Corollary~\ref{cor:Billey}~\ref{coro4.2_a}), we obtain the desired equality.
\end{proof}

The following lemma gives a recurrence of the equivariant $K$-theoretical 
Little\-wood--Richardson coefficients $c^w_{u,v}$.
We use the same idea as~\cite[Corollary~6.5]{M} and~\cite[Proposition~3.1]{PY}.

\begin{lemm}
\label{lem:recLR}
Let $u$, $v$, $w \in W^J$ and $s \in W^J$ a simple reflection.
If $c^w_{s,w} \neq c^u_{s,u}$, then we have
$$
c^w_{u,v}
 =
\frac{ 1 }
 { c^w_{s,w} - c^u_{s,u} }
\left(
 \sum_{u < x \le w} c^x_{s,u} c^w_{x,v}
 -
 \sum_{u \le y < w} c^w_{s,y} c^y_{u,v}
\right).
$$
In particular, we have
\begin{equation}
\label{eq:recLR}
c^w_{u,w}
 =
\frac{ 1 }
 { c^w_{s,w} - c^u_{s,u} }
\sum_{u < x \le w}
 c^x_{s,u} c^w_{x,w}.
\end{equation}
\end{lemm}

\begin{proof}
Consider the associativity
$$
\left( [\OO_s] [\OO_u] \right) [\OO_v]
 =
[\OO_s] \left( [\OO_u] [\OO_v] \right).
$$
Taking the coefficients of $[\OO_w]$ in the both hand sides and using Lemma~\ref{lem:LR}, 
we have
$$
c^u_{s,u} c^w_{u,v}
 +
\sum_{u < x \le w} c^x_{s,u} c^w_{x,v}
 =
c^w_{s,w} c^w_{u,v}
 +
\sum_{u \le y < w} c^w_{s,y} c^y_{u,v},
$$
from which we get the conclusion. 
\end{proof}

The Chevalley formula gives a combinatorial expression of $c^w_{s,v}$ for a simple reflection $s$.
To state the Chevalley formula of~\cite{LS} we need several notations.
For a dominant weight $\lambda \in \Lambda$, we put
$$
\mathbb{H}_\lambda
 =
\{ (\gamma^\vee, k) : 
 \gamma^\vee \in \Phi^\vee_+, \ k \in \Int, \ 0 \le k < \langle \gamma^\vee, \lambda \rangle
\}.
$$
Fix a total order on $I$ so that $I = \{ i_1 < \cdots < i_r \}$, 
and define a map $\iota : \mathbb{H}_\lambda \to \Rat^{r+1}$ by
$$
\iota \left(\gamma^\vee , k \right)
 =
\frac{ 1 }
 { \langle \gamma^\vee, \lambda \rangle }
\left( k, c_1, \ldots, c_r \right).
$$
where $c_1, \dots, c_r$ are the coefficients of the simple roots in $\gamma^\vee$ 
given by $\gamma^\vee = \sum_{j=1}^r c_j \alpha^\vee_{i_j}$.
Then it is known that $\iota$ is injective.
We define a total ordering $<$ on $\mathbb{H}_\lambda$ by
$$
h < h' \Longleftrightarrow \iota(h) <_{\lex} \iota(h'),
$$
where $<_{\lex}$ is the lexicographical ordering on $\Rat^{r+1}$.
For $h = (\gamma^\vee, k)$, we define affine transformations $r_h$ and $\tilde{r}_h$ on $\Lambda$ by
\begin{align*}
r_h(\mu) &= \mu - \langle \gamma^\vee, \mu \rangle \gamma,
\\
\tilde{r}_h(\mu) &= r_h(\mu) + \left( \langle \gamma^\vee, \lambda \rangle - k \right) \gamma.
\end{align*}
Note that $r_h = s_{\gamma}$.
Now we can state the Chevalley formula for the equivariant $K$-theory of the partial flag variety $\XX$.

\setbox\toto=\hbox {\cite[Theorem~4.8~(4.12) and~(4.13)]{LS}, see also~\cite[Corollary~7.1]{LP}}
\begin{prop}[\box\toto]
\label{prop:Chevalley}
%\textup{(\cite[Theorem~4.8 (4.12) and (4.13)]{LS}, see also~\cite[Corollary~7.1]{LP})}
Let $s$ be a simple reflection such that $s \in W^J$ and $v$, $w \in W^J$.
If $s = s_i$ and $\lambda = \lambda_i$ is the corresponding fundamental weight, then we have
\begin{equation}
\label{eq:Chevalley}
c^w_{s,v}
 =
\begin{cases}
 1 - e^{\lambda - v \lambda} &\text{if $w=v$},
\\
 \displaystyle\sum_{(h_1, \ldots, h_r)}
 (-1)^{r-1}
 e^{\lambda - v \tilde{r}_{h_1} \cdots \tilde{r}_{h_r} \lambda}
 &\text{if $w > v$,}
\\
 0 &\text{otherwise,}
\end{cases}
\end{equation}
where the summation is taken over all sequences $(h_1, \ldots, h_r)$ of length $r \ge 1$ 
satisfying the following two conditions:
\begin{enumerate}[(H1)]
\item\label{prop4.6_H1} %[(H1)]
$h_1 > h_2 > \cdots > h_r$ in $\mathbb{H}_\lambda$,
\item\label{prop4.6_H2} %[(H2)]
$v \lessdot v r_{h_1} \lessdot v r_{h_1} r_{h_2} \lessdot \cdots \lessdot v r_{h_1} \cdots r_{h_r} = w$ 
is a saturated chain in $W^J$.
\end{enumerate}
\end{prop}

\subsection{%
Connection to \texorpdfstring{$d$}{d}-complete posets
}\label{subsec4.3}

In this subsection we rephrase the Billey-type formula and the Chevalley-type formula 
in terms of combinatorics of $d$-complete posets.

Let $P$ be a connected $d$-complete poset with top tree $\Gamma$.
We regard $\Gamma$ as a simply-laced Dynkin diagram with node set $I$.
Let $\alpha_P$ and $\lambda_P$ be the simple root and the fundamental weight 
corresponding to the color $i_P$ of the maximum element of $P$.
We apply the results of Subsections~\ref{subsec4.1} and~\ref{subsec4.2} 
to the Kashiwara thick partial flag variety 
$\XX = \GG/\PP_-$, where $\PP_-$ is the maximal parabolic subgroup corresponding 
to $J = I \setminus \{ i_P \}$.
In this case, the parabolic subgroup $W_J$ coincides with the stabilizer $W_{\lambda_P}$ of $\lambda_P$ in $W$, 
and the minimum length coset representatives $W^J$ is denoted by $W^{\lambda_P}$.

By using a labeling of the elements of $P$ with $p_1, \ldots, p_N$ ($N= \# P$) 
so that $p_i < p_j$ in $P$ implies $i<j$, 
we can associate to each subset $D = \{ i_1, \dots, i_r \}$ ($i_1 < \dots < i_r$) of $P$ 
a well-defined element $w_D = s(p_{i_1}) \cdots s(p_{i_r}) \in W$.
Then the following formula is obtained from the Billey-type formula.

\begin{prop}
\label{prop:Billey_P}
Let $P$ be a connected $d$-complete poset and $F$ an order filter of $P$.
Then we have
\begin{equation}
\label{eq:Billey_P}
\xi^{w_F}|_{w_P}
 =
\sum_{E \in \EE^*_P(F)}
 (-1)^{\# E - \# F}
 \prod_{p \in E} \left( 1 - \vectz[H_P(p)] \right),
\end{equation}
under the identification $z_i = e^{\alpha_i}$ \textup{(}$i \in I$\textup{)}.
\end{prop}

\begin{proof}
Follows from Proposition~\ref{prop:Billey} 
by using Proposition~\ref{prop:root}~\ref{prop2.10_b} and Proposition~\ref{prop:E*=R*}.
\end{proof}

Also the following explicit expression is obtained from the Chevalley-type formula.

\begin{prop}
\label{prop:Chevalley_P}
Let $P$ be a connected $d$-complete poset and put $s = s_{i_P}$.
For two order filters $F$ and $F'$ of $P$, we have
\begin{equation}
\label{eq:Chevalley_P}
c^{w_{F'}}_{s,w_F}
 =
\begin{cases}
 1 - \vectz[F]
 &\text{if $F' = F$,}
\\
 (-1)^{\# (F' \setminus F) -1} \vectz[F]
 &\text{if $F' \supsetneq F$ and $F' \setminus F$ is an antichain,}
\\
 0
 &\text{otherwise,}
\end{cases}
\end{equation}
under the identification $z_i = e^{\alpha_i}$ \textup{(}$i \in I$\textup{)}.
\end{prop}

First we consider the case $r=1$ in Proposition~\ref{prop:Chevalley}.

\begin{lemm}
\label{lem:r=1}
Let $F$ be an order filter of $P$ and $h = (\gamma^\vee,k) \in \mathbb{H}_{\lambda_P}$.
If $w_F r_h \in W^{\lambda_P}$ and $w_F \lessdot w_F r_h \le w_P$, 
then there exists $p \in P$ such that $F' = F \sqcup \{ p \}$ is an order filter of $P$, 
$w_F r_h = w_{F'}$ and $\gamma^\vee = \gamma^\vee(p)$.
In this case $k=0$ and $\tilde{r}_h \lambda_P = \lambda_P$.
\end{lemm}

\begin{proof}
Since the interval $[e, w_P]$ in $W^{\lambda_P}$ is isomorphic to the poset of order filters of $P$ 
(Proposition~\ref{prop:Weyl}~\ref{prop2.11_a}), 
there exists a unique order filter $F'$ of $P$ such that $F' \supset F$, 
$\# F' = \# F + 1$ and $w_{F'} = w_F r_h$.
Hence we have $p \in P$ such that $F' = F \sqcup \{ p \}$ and $w_{F'} = s(p) w_F$.
We take a linear extension of $P$ such that $F = \{ p_{n+1}, \ldots, p_N \}$ 
with $N = \# P$ and $n = \# (P \setminus F)$.
If $p = p_m$, then $p$ is incomparable with $p_{m+1}, \ldots, p_n$, 
hence $s(p_m)$ is commutative with $s(p_{m+1}), \ldots, s(p_n)$ by Property~\ref{prop2.5_C5}
in Proposition~\ref{prop:dc_color}, 
and thus $s(p_i)\alpha^\vee(p_m)=\alpha^\vee(p_m)$ for $m+1 \le i \le n$.
Hence we have
\begin{align*}
\gamma^\vee
 &= 
w_F^{-1} \alpha^\vee(p_m)
\\
 &=
s(p_N) \cdots s(p_{n+1}) \alpha^\vee(p_m)
\\
 &=
s(p_N) \cdots s(p_{n+1}) s(p_n) \cdots s(p_{m+1}) \alpha^\vee(p_m)
\\
 &=
\gamma^\vee(p_m).
\end{align*}
By Proposition~\ref{prop:root}~\ref{prop2.10_c}, we see that $k=0$ and $\tilde{r}_h \lambda_P = \lambda_P$.
\end{proof}

Now we deduce Proposition~\ref{prop:Chevalley_P} from the Chevalley-type formula.

\begin{proof}[Proof of Proposition~\ref{prop:Chevalley_P}]
It follows from Proposition~\ref{prop:Weyl}~\ref{prop2.11_b}
and Proposition~\ref{prop:Chevalley}~that
$$
c^{w_F}_{s, w_F} = 1 - \vectz[F].
$$

Suppose that there exists a sequence $(h_1, \ldots, h_r)$ of elements in $\mathbb{H}_{\lambda_P}$ 
satisfying Conditions~\ref{prop4.6_H1} and~\ref{prop4.6_H2} in Proposition~\ref{prop:Chevalley}.
Then by Lemma~\ref{lem:r=1}, we have a sequence $(q_1, \ldots, q_r)$ of elements of $P$ such that 
$F_i = F \sqcup \{ q_1, \ldots, q_i \}$ is an order filter of $P$, $h_i = (\gamma^\vee(q_i),0)$ 
for $1 \le i \le r$ and $\tilde{r}_{h_1} \cdots \tilde{r}_{h_r} \lambda_P = \lambda_P$.
Now we show that $\{ q_1, \ldots, q_r \}$ is an antichain.
Assume to the contrary that there exist $i$ and $j$ such that $i<j$ and $q_i$ and $q_j$ are comparable.
Since $q_i$ is maximal in $P \setminus (F \sqcup \{ q_1, \ldots, q_{i-1} \})$ 
and $q_j \in P \setminus (F \sqcup \{ q_1, \ldots, q_{i-1} \})$, 
we have $q_i > q_j$.
Then by Proposition~\ref{prop:root}~\ref{prop2.10_a}, we see that $\gamma^\vee(q_i) < \gamma^\vee(q_j)$.
Hence by the definition of the total order on $\mathbb{H}_{\lambda_P}$, we have $h_i < h_j$,
which contradicts to Condition~\ref{prop4.6_H1}.
Moreover it follows from Proposition~\ref{prop:Weyl}~\ref{prop2.11_b}~that
$$
e^{\lambda_P - w_F \tilde{r}_{h_1} \cdots \tilde{r}_{h_r} \lambda_P}
 = \vectz[F'].
$$

Conversely, suppose that $F' \supsetneq F$ and $F' \setminus F$ is an antichain.
For $q \in F' \setminus F$, we put $h(q) = (\gamma^\vee(q),0) \in \mathbb{H}_{\lambda_P}$.
Since $F' \setminus F$ is an antichain, we can label the elements of $F' \setminus F$ 
so that $h(q_1) > \cdots > h(q_r)$.
Then $(h(q_1), \ldots, h(q_r))$ is the unique sequence satisfying Conditions~\ref{prop4.6_H1} and~\ref{prop4.6_H2} 
in Proposition~\ref{prop:Chevalley}.
\end{proof}


\section{%
Proof and corollaries of Main Theorem
}\label{sect5}

In this section, we give a proof of the main theorem (Theorem~\ref{thm:main} in Introduction) 
and derive several consequences.

\subsection{
Proof of the Main Theorem
}

Recall the main theorem of this paper:

\begin{theo}
\label{thm:main1}
Let $P$ be a $d$-complete poset and $F$ an order filter.
Then the multivariate generating function of $(P \setminus F)$-partitions, 
where $P \setminus F$ is viewed as an induced subposet of $P$, is given by
\begin{equation}
\label{eq:main1}
\sum_{\sigma \in \AP(P \setminus F)} \vectz^\sigma
 =
\sum_{D \in \EE_P(F)}
 \frac{ \prod_{v \in B(D)} \vectz[H_P(v)] }
 { \prod_{v \in P \setminus D} ( 1 - \vectz[H_P(v)] ) },
\end{equation}
where $D$ runs over all excited diagrams of $F$ in $P$.
\end{theo}
